\documentclass[11pt]{article}
\usepackage[margin=1in]{geometry}
\usepackage{amsmath,amssymb,amsthm}
\usepackage[colorlinks=true,urlcolor=blue,linkcolor=blue]{hyperref}
\newtheorem{theorem}{Theorem}
\newtheorem{lemma}[theorem]{Lemma}
\newcommand{\cost}{\operatorname{cost}}
\newcommand{\OPT}{\operatorname{OPT}}
\newcommand{\MS}{\mathrm{MS}}
\setlength{\emergencystretch}{4em}
\title{A Cha-Cha comparator for the exact multi-splay forest word}
\author{Research note; hand argument under audit}
\date{2 October 2026}
\begin{document}
\maketitle

\noindent\textbf{Status and dependencies.}
This note attempts the comparator transport left open in
\texttt{common-forest-normal-form.tex}. It uses the balanced-auxiliary
unfolding and exposed-splice lemmas of the local Cha-Cha manuscript.
It is a hand proof, not a Lean theorem. Its endpoint count is supplied
by \texttt{endpoint-embedding.tex}; its eventual implication for standard
multi-splay additionally depends on the conditional forest extension
of the current splay manuscript. In particular it does not prove the
exact $O(\log\log\log n)$ competitive ratio for standard multi-splay.

\section{The same keys and the same operation word}

Let $P$ be a perfect node-storing reference BST on $n=2^H-1$ keys.
View it as a routing tree on its $n+1$ external gaps. Its internal
routing keys are exactly the original keys, at depths $0,\ldots,H-1$.
This view introduces neither a second copy of a key nor additional
connector keys. In particular, no fixed maximum-key root from the
leaf-key reduction is added: the physical universe here is directly
the $n$ separators and their $n+1$ external gaps.

For a primary preferred path starting at depth $a$, retain the
depth-exit alphabet $[a,H]$. Its virtual routing labels are
$a,\ldots,H-1$; after unfolding its physical key set is precisely the
set of original nodes on that preferred path. All other primary paths
are separate components. Their positions in its search gaps can be
remembered analytically, but no such component is entered for free.

The common forest interface permits unit-cost selection of a named
component root, but no jump to an interior node. Each standard
multi-splay switch at $y$ has the exact three-split, three-join word
from \texttt{common-forest-normal-form.tex}. For a left-to-right switch,
write $U$ for the preferred prefix above $y$, $L$ for the old left
tail, $R$ for the incoming right tail, and $\ell,r$ for the closest
lower and upper ancestor keys in $U$. When both guards exist, the word is
\begin{align*}
 &\operatorname{Split}(y),\quad
   \operatorname{Split}(\ell),\quad
   \operatorname{Split}(r),\\
 &\operatorname{Join}(A,\ell,\varnothing),\quad
   \operatorname{Join}(R,r,B),\quad
   \operatorname{Join}(A',y,B').                 \tag{1}
\end{align*}
Here the second and third splits select the left and right daughters
of the first; the old tail $L$ is the second split's right daughter.
The omitted-guard and reflected cases use the corresponding shortened
word. The word is determined by $P$, the initial preferences, and the
original request sequence, independently of all auxiliary tree shapes.

Process this word in its actual order: wrong proper ancestors in
decreasing depth, followed by the final switch at the requested node.
For each switch at depth $d$, the comparator selects its named path
component and performs one local depth request to $d$. It does not
perform a depth-$H$ completion. The final split at the requested node
itself serves that original key. At a primary leaf the final switch
may be represented by toggling the analytical preference between its
two empty external gaps. This changes no actual auxiliary membership,
but its guard splays and its final depth emission are retained.

Let $S$ count all these switches, including the mandatory final
switches. Let $F_{\mathrm{act}}$ be the all-right-initialized fixed-owner
secondary variation of exactly this emission stream, in the actual
order, on a fixed balanced interval tree $B$ over $0,\ldots,H$. Put
\[
 K=1+\operatorname{height}(B)=O(1+\log(H+2)),\qquad
 \lambda=1+\log_2 K.
\]

\section{Local suffix accesses and the histories they really write}

The representation of a path starting at $a$ consists of genuine
canonical intervals $I=[s,t]$ with $s\ge a$, together with a structural
boundary spine $\beta_a$. The latter consists of surviving ancestors
of leaf $a-1$ with split $c\ge a$. It is an all-left chain in the
induced suffix tree, stored in its own balanced auxiliary. Genuine
canonical subtrees hang in its gaps. The seam between its final left
child and the leftmost genuine canonical subtree is administrative:
these are stored in separate auxiliaries even though the boundary
direction is left. Every auxiliary has at most $K$ keys.

\begin{lemma}[A local suffix transaction]\label{lem:suffix}
A depth request $d\ge a$, followed by a primary splice at depth $d$,
can be performed on this suffix representation in
$O((1+f)\lambda)$ pointer and rotation cost, where $f$ counts changes
only at genuine secondary intervals visited by this request. It restores
the same suffix convention at the unchanged starting depth $a$ and
preserves all deeper-owner histories. The old and incoming tails use
the suffix convention at $d+1$.
\end{lemma}

\begin{proof}
First concatenate $\beta_a$ with the root auxiliary of the genuine
subtree at its final left gap, if both exist. This removes the one
administrative seam and gives an ordinary Tango representation of the
induced suffix reference, with all boundary preferences left. This
concatenated auxiliary is a reference path and still has at most $K$
keys. One balanced split/join interface costs $O(\lambda)$.

Apply the ordinary Tango access to $d$ in this induced reference. Its
search path leaves the boundary chain through a right edge at most
once. Thus at most one boundary preference changes, in addition to
the $f$ genuine changes. The access costs $O((1+f)\lambda)$ and exposes
the entire induced reference path to leaf $d$ in the root auxiliary.

Apply the exposed splice to this root auxiliary. Its labels greater
than $d$ are exactly the structural boundary of the suffix $[d+1,H]$;
the usual proof by surviving ancestors of leaf $d$ is unchanged by
the restriction to $[a,H]$. Split them off, expose the final two prefix
exit gaps, exchange their roles, and concatenate the incoming
$\beta_{d+1}$. These are constantly many balanced auxiliary operations,
of total cost $O(\lambda)$.

Every crossing genuine interval has owner depth $s\ge a$, contains
$d$, and was just written toward $d$. Its preference therefore agrees
with the new exposed path. Intervals with $s<a$ are structural boundary
nodes, not genuine histories. Intervals with $s>d$ travel with their
respective tails and retain their bits.

Finally restore the unique boundary preference that may have changed
to right. Its split $c$ satisfies $c<d$, so this boundary node survived
in the prefix and its old left continuation was not affected by the
primary splice. One ordinary Tango switch makes it left again. All
boundary preferences are now left and their full chain belongs to the
root auxiliary. Separate its final genuine continuation at the
administrative seam. These two final operations cost $O(\lambda)$,
do not write any genuine bit, and restore the required representation.
If a spine, continuation, or tail is empty the corresponding operation
is omitted.
\end{proof}

\begin{lemma}[Deleted writes, not fictional history updates]\label{lem:deletion}
Across all actual switches, the number $G$ of genuine secondary
changes made by these local suffix accesses satisfies
\[
 G\le F_{\mathrm{act}}.
\]
\end{lemma}

\begin{proof}
Fix one owner row $(I,u)$, where $I=[s,t]$. Delete from its full
$F_{\mathrm{act}}$ write word every event whose selected primary path
starts strictly below $s$. All remaining writes are exactly the
comparator's genuine writes in this row.

Indeed, on a switch at depth $d$ the current preferred path agrees
with the original requested node between its starting depth $a$ and
$d$. Thus whenever $s\ge a$ and $I$ contains $d$, its actual owner is
the fixed primary ancestor at depth $s$ of that request. Whenever
$s<a$, the surviving reference node is structural and supplies no
independent history. A genuine owner bit exists in the component
containing its owner, travels with that owner, and changes only on
these genuine accesses. Lemma~\ref{lem:suffix} shows that a splice
neither resets it nor introduces an unrecorded genuine write.

Both words start from the same initial right bit. Deleting entries
from a binary word cannot increase its variation, by the triangle
inequality between the surviving consecutive bits. Apply this fact
to each fixed owner row and sum. The deletion rule may depend on the
preceding primary states; variation monotonicity needs no independence.
\end{proof}

\section{Exposing the original node guards}

The nontrivial transport issue is that a virtual depth label and its
physical unfolding image need not be the same primary separator.
Consequently the presence of routing labels $d-1,d$ in the root
auxiliary alone does not identify the standard multi-splay guards.
The following argument uses actual external gaps instead.

\begin{lemma}[Guard exposure with the same node alphabet]\label{lem:guards}
After the local inner access to depth $d$, an additional $O(\lambda)$
rooted preparation exposes a physical hull containing the switched
key $y$ and both existing nearest-ancestor guards $\ell,r$. This can
be done without changing the common forest domains. It holds even
when one or both primary tails have no internal keys.
\end{lemma}

\begin{proof}
For notation consider a left-to-right switch. In the virtual depth
tree, cut after routing label $d$ as in the exposed-splice lemma,
retaining $d$ with the prefix. Treat the old tail as one opaque
terminal hole. This is a representation of the same current component
$U\cup\{y\}\cup L$: its keys in $L$ remain attached at that hole.
The preparation consists of rotations inside the current component;
it does not execute an early common-forest split.

The last two prefix depth leaves are the off-path exit $E_d$ and the
terminal hole at $d+1$. Routing keys $d$ and, if $d>a$, $d-1$ lie in
the exposed balanced auxiliary. Their two prefix search paths therefore
have $O(\lambda)$ nodes. Rooted-hull unfolding gives physical paths
of that same order. Explicitly walk both paths in the prefix skeleton,
stopping at their gap or at the first root of the opaque attached
subtree. This statement includes null external gaps; it does not
depend on the presence of a header to read.

In the actual prefix alphabet $U\cup\{y\}$, the old left-tail gap
has neighboring keys $\ell,y$, omitting $\ell$ if it does not exist.
The incoming right-tail gap has neighbors $y,r$, with the analogous
omission. Both neighbors of an external BST gap lie on its search path:
the lower neighbor is the last lower bound and the upper neighbor the
last upper bound. Consequently the two explicitly traversed physical
paths contain $y,\ell,r$ whenever present. The same argument applies
when the tails are empty, and reflection handles the opposite switch.

All rotations in the virtual preparation have a rooted physical hull
of size $O(\lambda)$; the two walks add only $O(\lambda)$ nodes.
The rooted touched-hull argument pulls this set back to a rooted hull
in the original component with the same keys and frontier subtrees.
At no point has the old tail ceased to belong to that component.
\end{proof}

\section{Inserting the prescribed split and join checkpoints}

\begin{lemma}[Checkpoint insertion]\label{lem:checkpoints}
The suffix transaction of Lemma~\ref{lem:suffix} admits a comparator
execution of precisely the word (1), or its omitted-guard or reflected
version, with cost $O((1+f)\lambda)$. Between consecutive operations
every rotation takes place inside a current common component.
\end{lemma}

\begin{proof}
Carry out the local access and guard preparation inside the selected
component. The incoming tail is selected through its own root handle.
The balanced boundary concatenations may inspect $O(\lambda)$ nodes
of that second component. Keep its inspected set as a second rooted
hull, rather than treating this root selection as an interior jump.
The common interface permits such paid selections and preparation of
another live component between prescribed operations; they do not
change its domain. Equivalently, any preparation entirely inside the
incoming component can be postponed until immediately before the join
that consumes it.

The exposed-splice construction supplies a bounded frontier certificate:
before and after the splice, all untouched physical frontier subtrees
are the same ordered key intervals; the union of touched hull keys
has size $O((1+f)\lambda)$. The old and incoming tail roots are included
when nonempty. This also follows by the usual touched-skeleton tracking:
a move or rotation exposes only its new gap root, and detaching or
attaching a tail exposes its root. Thus each resulting component has
a rooted after-skeleton and the unchanged frontier components.

Use the guard-containing skeleton to prepare the original component
in the form
\[
 y\bigl(\ell(A,L),\ r(\varnothing,B)\bigr).
                                                        \tag{2}
\]
Only skeleton keys are rotated. No opaque frontier subtree straddles
one of $y,\ell,r$, since all these keys are in the skeleton. Hence
they can all be preserved while the skeleton is put in (2), in time
linear in its size. The subtrees displayed in (2) include all their
opaque hanging components. The desired internal skeleton shapes
inside $A,L,B$ can be arbitrary at this stage.

Now execute the three specified splits. Their pivots are already
component roots, so each is a legal common root deletion. The old
tail $L$ becomes a separate component exactly at its specified split,
and no earlier. Perform the three raw root joins of (1), with the
prescribed domains and pivots. They give
\[
 y\bigl(\ell(A,\varnothing),\ r(R,B)\bigr)
                                                        \tag{3}
\]
and the parked component $L$.

Finally reconfigure the portions of the skeleton in these resulting
components to the valid after-representations prescribed by the
exposed splice and by boundary restoration. Each is a rooted skeleton
with the same ordered frontier intervals; linear-cost BST reconfiguration
therefore applies separately inside each component. No rotation crosses
a component boundary. An absent guard simply removes its split, join,
and displayed skeleton node. If both tails are empty, the same
guard preparation and checkpoints are performed; the argument never
relied on a nonempty tail.

The preparation, six checkpoints, and final reconfigurations cost
$O((1+f)\lambda)$ in total. They end with the required persistent owner
histories, because those histories were specified by the ideal local
suffix transaction and have not been initialized again.
\end{proof}

\begin{theorem}[The proposed common-word comparator]\label{thm:comparator}
The exact multi-splay forest word has a comparator of cost
\[
 C_{\mathrm{forest}}(W_{\MS}(X))
   =O\bigl(n+\lambda(S+F_{\mathrm{act}})\bigr).
                                                        \tag{4}
\]
The endpoint count theorem consequently gives
\[
 C_{\mathrm{forest}}(W_{\MS}(X))
   =O\bigl(\lambda(\OPT(X)+n)\bigr).
                                                        \tag{5}
\]
\end{theorem}

\begin{proof}
Initialize every path's balanced suffix representation with the
all-right genuine histories, in total $O(n)$ work. A component with
$k$ keys can be changed to its prescribed initial shape in $O(k)$
rooted rotations and moves. Sum Lemma~\ref{lem:checkpoints} over the
$S$ switches, and use Lemma~\ref{lem:deletion}. This gives (4) with
no lifetime initialization charge and no persistent splay potential:
the comparator's auxiliaries are balanced.

The separately audited endpoint embedding and actual-order lemma give
$S+F_{\mathrm{act}}=O(\OPT(X)+n)$ for this exact emission stream.
Substitution yields (5).
\end{proof}

\section{Padding the reference without adding physical keys}

The comparator construction also applies to the customary minimum-height
balanced reference on an arbitrary number $n$ of keys. Put
$H=\lceil\log_2(n+1)\rceil$. Complete its routing reference to depth
$H$ by adding conceptual dummy routing nodes below each original
external gap. Within each such padded gap region, retain only its
all-right depth-$H$ descendant as an \emph{active} gap representative.
Every other padded leaf is inactive. There are still exactly $n+1$
active external gaps, in their original order.

Give every dummy primary vertex the fixed right preference. It has
at most one active child, whereas every original vertex has two active
children. Suppressing the inactive pieces and unary refinements
therefore yields exactly the original node-storing reference, with
each original separator once and no physical dummy key. Complete the
secondary virtual alphabet to $0,\ldots,H$ as before.

The unfolding proof continues to hold with empty components suppressed:
a depth interval has at most two active interval components, an
abstract routing node has at most two physical images, and the same
rooted-hull and equal-frontier proofs apply. Here an active external
gap is retained as an ordered atom even when it contains no stored
key. Thus the two null child gaps at a real primary leaf remain
different active atoms separated by that leaf's real key. In
particular the two-gap guard argument does not disappear when both
tails are physically empty.

Every emitted depth belongs to a real switched node. Every owner
ancestor at an earlier or equal depth is therefore real as well.
No genuine row whose owner is dummy ever receives a write; all such
bits stay right. A physically empty tail lies entirely in a dummy
region, so all its genuine histories and its structural boundary can
be reconstructed from its starting depth and region. Its auxiliary
shape can be chosen canonically without physical work. This introduces
no potential reset, because the comparator uses balanced auxiliaries
and has no shape potential. A nonempty component stores all its
nonconstant virtual histories in the same $O(H)$-bit header as before.

Consequently the proof of (4) holds unchanged for arbitrary $n$ and
the exact original forest word. No split or join at a dummy pivot is
inserted. The completed reference has fewer than $2n+1$ internal
vertices, so any linear reference-count initialization allowance also
remains $O(n)$. The active-padding endpoint corollary in
\texttt{endpoint-embedding.tex} supplies (5) for this reference too.

Every intermediate component in the word has at most $H$ keys.
Combining (5) with the conditional common-forest splay comparison
would give
\[
 \cost_{\MS}(X)=
 O\!\left((\OPT(X)+n)
     (1+\log\log H)^2
     (1+\log\log\log H)^2\right).
\]
For $H=\Theta(\log n)$ this is the weaker factor
$O((1+\log\log\log n)^2
    (1+\log\log\log\log n)^2)$.
The local comparator proof and the forest extension require their
own audit; neither the known Cha-Cha theorem nor the endpoint
lower bound alone establishes this conclusion.

\end{document}
